% !TeX root = ../papers/nc-p3.tex
\section{階数14の二位置移動による降下}
\label{sec:rank-fourteen-descent}
\begin{proposition}\label{prop:rank-fourteen-descent}
定理\ref{thm:fixed-opposite-rank-drop}の階数14の二成分は、それぞれ通常倍率の全域で
二位置移動により有理な $t$ 直線の開集合に降りる。
各成分の中心の階数と四中心による有限商の検査は、一変数の局所化多項式環上の行列で実行できる。
\end{proposition}
\begin{proof}
$i^2=-1$ を固定し、前定理の $w=\lambda r$ に対し $s^2=w$ と置く。
第一行列の第一行を $(rX,b)$ とする。全域で使える表示は
\[
 X(s)=i\big(s^2-(1+i)s-1\big),\qquad b(s)=s^2+(1-i)s-1 .
\]
この対は $K(X,b)=0$ を恒等的に満たす。$X/i-b=-2s$ で、$s=0$ でも同時には消えない。
$(E,F)=((1-i)b+iX,(1+i)X-ib)$ を代入すると
$(E(s),F(s))=(-X(-s),-b(-s))$ となる。従って $B(s)=-A(-s)$ である。
二位置移動は二つの行列を交換するため、共通スカラーを除き $s\mapsto-s$ と一致する。
この二次拡大の固定体は $\mathbb Q(i)(t)$ であり、各成分の二点のAS性は等しい。

通常倍率を与える $t$ の開部分の座標環を $R$ とする。$r=0,\infty$ は非通常なので
この開部分はアフィン直線に含まれ、$R$ は $\mathbb Q(i)[t]$ の局所化である。
$w$ は単元で、$S=R[s]/(s^2-w)$ は階数2の有限エタール代数となる。
上の表示の行列式は、単元 $r^2$ を除き
\[
 X^2-b^2=-2(s^2-1)(s-i)^2,\qquad E^2-F^2=-2(s^2-1)(s+i)^2
\]
である。いずれの因子も通常倍率で消えず、関係空間を表す元のテンソルを $S$ 上で構成できる。
従って $b=0$ を別の座標で張り替える必要もない。$c+ds$ による乗算は $R$ 基底 $(1,s)$ では
\[
 \begin{pmatrix}c&wd\\d&c\end{pmatrix}
\]
である。これで経路空間の関係の各挿入、中心の交換方程式、選んだ中心式の挿入を
$R$ 上の行列に移せる。道空間の商を作るときも、表示から核や余核を取ればよい。
$R$ は単項イデアル整域なので、標準次元が既知の長さ4,5の道空間は
$R$ 加群として有限自由である。階数条件は小行列式の最大公約数で判定できる。
幾何ファイバーで制限した行列は二つの $s$ に対応する行列の直和になり、
全始点での検査は二位置移動で同値になる。
\end{proof}

\begin{proposition}\label{prop:rank-fourteen-residue}
階数14のこの表示を標数零の離散付値環上で取り、剰余体を $\mathbb F_{101}$ とする。
その還元が通常な $\mathbb F_{101}$ 有理点ならば、対応する分数体上の窓は
標準Koszul左右AS正則である。該当する剰余点は $i$ の二値で各96点ある。
\end{proposition}
\begin{proof}
101では $i=10,91$ である。\path{computations/verify_rank_fourteen_descent.py} は
全ての $t$ と $s^2=w(t)$ の全根を列挙する。$t=\infty$ では $r=0$ となり通常性に反する。
許容な192点の全てで隣接階数6、反対階数14、経路長4,5の次元35,56と四次中心の文字次元
$(3,2,0,0)$ を検算する。最初の二文字から二式ずつ選ぶ。
四つの剰余点では文字0の二つの線型結合を取り直すと、
全192点・全四始点で商次元が命題\ref{prop:fixed-residual-regularity}の表に一致する。
同命題の付値環上の持上げにより、各剰余点に還元される窓で全次数の結論を得る。
\end{proof}
